% Canon insertion for SST_CANON-v0.8.35.tex
% Recommended location: immediately after Subsection "Circulation Quantization"
% and before "Bounded-Domain Biot--Savart Admissibility".

\subsection{Material Swirl--Tonic Potential and Vorticity Flux}
    \label{subsec:material_swirl_tonic_potential}

    \textbf{[ORTHODOX KINEMATIC IDENTITY / CANONICAL NOTATION].}
    On every smooth material Euler branch, SST may use the material velocity itself
    as a vorticity-potential representation.  Define the \emph{material swirl--tonic
    potential}
    \begin{align}
        \mathbf A_{\rm st}^{(m)}
        &:= \mathbf v,
        &
        \boldsymbol{\zeta}
        &:= \nabla\times\mathbf A_{\rm st}^{(m)}
        =\nabla\times\mathbf v.
        \label{eq:material_swirl_tonic_definition}
    \end{align}
    The superscript ``$(m)$'' denotes the material Euler sector.  It is retained to
    prevent confusion with the independent transverse/link field
    $\mathbf A_{\mathrm{eff}}$ of Sec.~\ref{sec:embridge}.

    In differential-form notation, with $\flat$ the Euclidean musical map,
    \begin{align}
        a_{\rm st}^{(m)}
        &:=(\mathbf v)^\flat,
        &
        \Omega_{\rm st}^{(m)}
        &:=d a_{\rm st}^{(m)}.
        \label{eq:material_swirl_tonic_forms}
    \end{align}
    The two-form $\Omega_{\rm st}^{(m)}$ is the vorticity-flux form associated with
    $\boldsymbol{\zeta}$.  Therefore, for every smooth orientable surface $S$ with
    boundary $C=\partial S$ on which the fields are regular, Stokes' theorem gives
    \begin{align}
        \boxed{
        \Gamma_C
        =
        \oint_C\mathbf A_{\rm st}^{(m)}\cdot d\boldsymbol\ell
        =
        \oint_C\mathbf v\cdot d\boldsymbol\ell
        =
        \int_S\boldsymbol{\zeta}\cdot d\mathbf S .
        }
        \label{eq:material_swirl_tonic_stokes}
    \end{align}
    This is an orthodox identity of smooth vector calculus and ideal-fluid
    kinematics \cite{MajdaBertozzi2002}.  Maxwell's 1855--1856 use of
    incompressible ``tubes of fluid motion'' and his later electro-tonic line/surface
    integral construction provide historical motivation for the terminology only;
    Maxwell explicitly treated the fluid as a mathematical representation rather
    than as an asserted material substrate \cite{Maxwell1856FaradayLines}.

    \paragraph{Dimensional closure.}
    The material swirl--tonic quantities satisfy
    \begin{align}
        [\mathbf A_{\rm st}^{(m)}]
        &=[\mathbf v]=\mathrm{m\,s^{-1}},
        &
        [\boldsymbol{\zeta}]
        &=\mathrm{s^{-1}},
        &
        [\Gamma_C]
        &=\mathrm{m^2\,s^{-1}}.
        \label{eq:material_swirl_tonic_dimensions}
    \end{align}
    Hence Eq.~\eqref{eq:material_swirl_tonic_stokes} is dimensionally closed without
    importing electromagnetic SI units.

    \paragraph{Normalized material holonomy.}
    Relative to the calibrated circulation scale of Eq.~\eqref{eq:gamma0}, define
    \begin{align}
        h_C^{(m)}
        :=
        \frac{1}{\Gamma_0}
        \oint_C\mathbf v\cdot d\boldsymbol\ell .
        \label{eq:material_swirl_tonic_holonomy}
    \end{align}
    In a multiply connected exterior domain, the collection of such periods probes
    the harmonic circulation sector described in
    Subsection~\ref{subsec:hodge_topological_circulation}.  A period labels a global
    circulation channel; by itself it does not identify knot species, chirality, or
    particle identity.

    \paragraph{No-gauge-redundancy guard.}
    The phrase ``vector potential'' in
    Eq.~\eqref{eq:material_swirl_tonic_definition} refers to the differential
    relation $\boldsymbol{\zeta}=\nabla\times\mathbf v$.  It does \emph{not} make the
    physical material velocity a gauge-redundant variable.  A replacement
    \begin{align}
        \mathbf v\longrightarrow\mathbf v+\nabla\chi
        \label{eq:material_swirl_tonic_nogauge}
    \end{align}
    leaves the local curl unchanged for smooth single-valued $\chi$, but generally
    changes material trajectories, kinetic energy, and boundary data.  Conversely,
    a harmonic-knot field can change non-contractible circulation periods without
    changing the local curl.  The physical material state is therefore fixed by the
    velocity together with its boundary conditions and circulation sector, not by
    vorticity alone \cite{CantarellaDeTurckGluck2002Hodge}.

    \paragraph{Material--link sector separation.}
    The Canon retains
    \begin{align}
        \boxed{
        \mathbf A_{\rm st}^{(m)}=\mathbf v
        \;\not\equiv\;
        \mathbf A_{\mathrm{eff}},
        \qquad
        \boldsymbol{\zeta}
        \;\not\equiv\;
        \nabla\times\mathbf A_{\mathrm{eff}}
        }
        \label{eq:material_swirl_tonic_link_guard}
    \end{align}
    unless a separate dynamical matching theorem is supplied.  The SI-normalized
    bridge $\mathbf A_{\rm SI}=(m_e/e)\,\mathbf u_{\!\boldsymbol{\circlearrowleft}}$
    already recorded in
    Eq.~\eqref{eq:canonical_em_prefactor_guard} is a declared normalization bridge;
    it does not collapse the material Euler and photon/link sectors into one field.

    \paragraph{Canonical falsification gate.}
    Any numerical calculation that uses
    Eq.~\eqref{eq:material_swirl_tonic_definition} must independently converge the
    line circulation and vorticity-flux evaluations,
    \begin{align}
        \epsilon_{\rm Stokes}
        :=
        \frac{
        \left|
        \oint_C\mathbf v\cdot d\boldsymbol\ell
        -
        \int_S\boldsymbol{\zeta}\cdot d\mathbf S
        \right|
        }{
        \max\!\left(
        \left|\oint_C\mathbf v\cdot d\boldsymbol\ell\right|,
        \left|\int_S\boldsymbol{\zeta}\cdot d\mathbf S\right|,
        \Gamma_{\rm num}
        \right)
        }
        \longrightarrow0
        \label{eq:material_swirl_tonic_stokes_falsifier}
    \end{align}
    under quadrature and resolution refinement, where $\Gamma_{\rm num}>0$ is a
    preregistered numerical floor used only to regularize a near-zero denominator.
    Persistent non-convergence falsifies the resolved field realization or its
    claimed regularity assumptions.  It does not license replacing the material
    field by $\mathbf A_{\mathrm{eff}}$.
